\documentclass[11pt]{article}
\usepackage[a4paper,margin=28mm]{geometry}
\usepackage{amsmath,amssymb,amsthm,hyperref}
\newtheorem{proposition}{Proposition}
\newtheorem{lemma}[proposition]{Lemma}
\newcommand{\Alt}{\operatorname{Alt}}
\newcommand{\Sym}{\operatorname{Sym}}
\title{A structural label in the subgroup table of $A_{11}$}
\author{human\_compression project}
\date{}
\begin{document}
\maketitle

Connor and Leemans' companion table for $A_{11}$, dated 27 November
2012, prints $2\cdot A_5:2$ in row 186, of order $240$ and class length
$20790$ \cite{table}. The native page is page 11. Section 3.2 of their
article distinguishes the nonsplit middle dot from the general extension
symbol and explicitly adopts left association \cite[p.~263]{atlas}.
Thus the printed expression denotes $(2\cdot A_5):2$ under the stated
convention. We prove that this interpretation is impossible in $A_{11}$.
We then identify the actual subgroup supplied by the independently checked
numerical table. The original expression remains in the source ledger;
the corrected structure is recorded separately.

\begin{lemma}
A nonsplit extension of $A_5$ by a normal subgroup of order two is perfect.
\end{lemma}
\begin{proof}
Write $1\to Z\to N\xrightarrow{q}A_5\to1$, with $|Z|=2$.
The subgroup $Z$ is central, since its automorphism group is trivial.
Recall that $A_5$ is perfect. For completeness, the commutator of
$(12)(34)$ and $(12)(35)$ is a $3$-cycle. The commutator subgroup
therefore contains every $3$-cycle, since these form one conjugacy class
in $A_5$; they generate $A_5$.
Consequently $q(N')=A_5$. If $N'\cap Z=1$, then $N'$ is a complement
to $Z$, contradicting the nonsplitting assumption. Hence $Z\subseteq N'$
and $N'=N$.
\end{proof}

\begin{proposition}
No nonsplit extension $2\cdot A_5$ embeds in $A_{11}$. In particular,
no subgroup of $A_{11}$ has the structure $(2\cdot A_5):2$.
\end{proposition}
\begin{proof}
Suppose that such a subgroup $N$ exists, and let $z$ be its central
involution. Since $z$ is an even permutation on eleven points, it is a
product of $m=2$ or $m=4$ disjoint transpositions. Put
$\Omega=\operatorname{supp}(z)$. Centrality makes $\Omega$ invariant
under $N$, and restriction defines a homomorphism
\[
 N\longrightarrow C_{\Sym(\Omega)}(z)
       \cong C_2^{\,m}\rtimes S_m.
\]
Indeed, a permutation commuting with $z$ permutes its two-point orbits
and independently interchanges the two points in each orbit.
The group on the right is solvable: its base group is abelian, and
$S_2$ and $S_4$ are solvable. One may use the normal series
$1\triangleleft V_4\triangleleft A_4\triangleleft S_4$ to see the latter.
The image of the perfect group $N$ is perfect, while every perfect
solvable group is trivial. But $z$ has nontrivial restriction to
$\Omega$, a contradiction. The final assertion follows because the
indicated semidirect product contains $2\cdot A_5$ as a subgroup.
\end{proof}

This obstruction concerns every subgroup of $A_{11}$; it does not depend
on the choice of a correspondence between the numerical source rows and
the computed representatives. It does not assert that the authors intended
the impossible interpretation. A missing pair of parentheses would change
the expression to a valid one.

\begin{proposition}
The saved representative corresponding to row 186 is
\[
 H=\langle (12345),\ (12)(8\,10\,9\,11),\ (8\,9)(10\,11)\rangle.
\]
It is isomorphic to
\[
 \{(\sigma,k)\in S_5\times C_4:
                   \operatorname{sgn}(\sigma)=(-1)^k\}.
\]
Its structure can be written $(A_5\times2)\cdot2$, or equivalently
$2\cdot(A_5:2)$, with the displayed parentheses.
\end{proposition}
\begin{proof}
Put $r=(8\,10\,9\,11)$, acting on a set disjoint from the first five
points. Each displayed generator belongs to the indicated pullback.
The first two generators project onto $S_5$: a $5$-cycle and an adjacent
transposition generate $S_5$. Moreover, the square of the second generator
is $r^2\ne1$. The projection kernel therefore has order at least two.
The pullback has order $2|S_5|=240$, so containment and projection show
that $H$ is the entire pullback.

The elements with even $k$ form $N=A_5\times\langle r^2\rangle$,
of index two. Every element outside $N$ has odd $k$, so its square has
second coordinate two and cannot be the identity. Thus the extension
of $C_2$ by $N$ has no complement and is nonsplit. This gives the first
expression. Projection onto $S_5$ has central kernel $\langle r^2\rangle$.
Every lift of a transposition has square $r^2$; a section of this
projection would preserve the order of a transposition, which is
impossible. Hence this central extension is nonsplit as well. Finally,
$S_5=A_5:C_2$, with the usual nontrivial conjugation action, giving the
second expression.
\end{proof}

The accompanying literal certificate uses only permutations on eleven
points. It closes the saved generators and checks equality with all $240$
pullback elements. It verifies the direct product of order $120$, the
absence of involutions outside it, and all twenty lifts of the ten
transpositions. As an additional check, it propagates all eight possible
assignments of the three generators to $C_2$ across every Cayley edge.
Exactly one assignment is surjective, and its kernel is the displayed
$A_5\times C_2$. The centre of that kernel has order two. Thus the
index-two subgroup cannot itself be a nonsplit extension $2\cdot A_5$.

The table and its naming convention are due to Connor and Leemans.
The arguments above use elementary facts about central extensions,
perfect groups, and permutation centralizers. The contribution here is
the explicit correction and its connection to the checked representative;
no general priority claim is made for those facts. This note corrects
one structural entry. It neither changes the numerical subgroup lattice
nor proves an open conjecture.

\begin{thebibliography}{9}
\bibitem{atlas} T. Connor and D. Leemans,
\emph{An atlas of subgroup lattices of finite almost simple groups},
Ars Mathematica Contemporanea \textbf{8} (2015), 259--266.
\bibitem{table} T. Connor and D. Leemans,
\emph{The subgroup lattice of $A_{11}$}, companion table,
27 November 2012, 41 pages; row 186 on page 11.
\end{thebibliography}
\end{document}
